\documentclass[11pt]{amsart}
\usepackage{geometry}                % See geometry.pdf to learn the layout options. There are lots.
\geometry{letterpaper}                   % ... or a4paper or a5paper or ... 
%\geometry{landscape}                % Activate for for rotated page geometry
%\usepackage[parfill]{parskip}    % Activate to begin paragraphs with an empty line rather than an indent
\usepackage{graphicx}
\usepackage{amssymb}
\usepackage{epstopdf}
\DeclareGraphicsRule{.tif}{png}{.png}{`convert #1 `dirname #1`/`basename #1 .tif`.png}

\title{PS 3.40}

%\date{}                                           % Activate to display a given date or no date

\begin{document}
\maketitle
%\section{}
%\subsection{}
\noindent Se tiene un grupo de registros donde se almacena información relativa a los ele­mentos de la tabla periódica de química. Por cada elemento se ingresa su nom­bre, su conductividad eléctrica y su conductividad térmica. Haga un diagrama de flujo para calcular lo siguiente :
\newline

\noindent a) Los dos principales elementos conductores de la electricidad y el calor.
\noindent b) Los dos peores elementos conductores de la electricidad y el calor.
\newline

\noindent Datos: $NOM_{1}$, $CE_{1}$, $CT_{1}$, $NOM_{2}$, $CE_{2}$,$CT_{2}$,...,"NN",0,0
\newline

\noindent Donde: 
\newline

\noindent NOM¡ es una variable de tipo cadena de caracteres que representa el nombre del elemento i.

\noindent CE¡ es una variable de tipo real que representa la conductividad eléc­trica del elemento i. Si no se conoce el valor del elemento i se in­gresa como CE el valor de -1.

\noindent CTi es una variable de tipo real que representa la conductividad térmi­ca del elemento i. Si no se conoce el valor del elemento i se ingresa como CT el valor de -1.
\newline

\noindent El fin de datos está expresado como : “NN”,0,0.

\end{document}  